La calibration vise à estimer expérimentalement les paramètres de bruit des
capteurs, afin de les modéliser correctement dans le filtre de Kalman Étendu.
Deux capteurs nécessitent une calibration : le module GNSS et l'IMU.

\subsubsection{Calibration du module GNSS}

La calibration du GPS consiste à estimer la matrice de bruit de mesure $R$,
qui représente l'incertitude des positions fournies par le récepteur. Cette
calibration a été réalisée lors du choix de la constellation satellitaire :
une acquisition statique extérieure a permis de mesurer les écarts-types
$\sigma_{\text{lon}}$ et $\sigma_{\text{lat}}$ pour chaque configuration.

\textbf{Remarque importante :} la calibration a été effectuée exclusivement
en extérieur, à l'écart de tout obstacle. Une acquisition en intérieur serait
inexploitable en raison du phénomène de \textit{multipath} : les signaux
satellitaires réfléchis par les murs et les plafonds produisent une
distribution bimodale des positions, incompatible avec l'hypothèse gaussienne
du filtre.

Les valeurs obtenues et la matrice $R$ résultante sont détaillées en
section~\ref{sec:params}.

\subsubsection{Calibration de l'IMU}

L'IMU nécessite deux calibrations distinctes : l'estimation du biais du
gyroscope, et l'estimation de l'écart-type de l'accélération $\sigma_a$
utilisé dans la matrice de bruit de processus $Q$.

\paragraph{Pour le biais du gyroscope}

Un gyroscope parfait indiquerait une vitesse angulaire nulle lorsque le
capteur est immobile. En pratique, les MEMS présentent un \textbf{biais
additif} $b_\omega$, constant à court terme, qui s'intègre au cours du temps
et produit une dérive du cap estimé :

\begin{equation}
    \omega_{\text{mesuré}} = \omega_{\text{réel}} + b_\omega + \epsilon
\end{equation}

où $\epsilon$ est le bruit de mesure centré. Non corrigé, ce biais entraîne
une erreur d'orientation croissant linéairement avec le temps, ce qui fausse
l'étape de prédiction de l'EKF.

\textbf{Protocole de calibration :} au démarrage de chaque acquisition, le
système est maintenu immobile pendant les premières secondes. Durant cette
phase quasi-statique, la vitesse angulaire réelle est nulle ($\omega_{\text{réel}} = 0$).
Le biais est estimé par la moyenne des $N = 200$ premières mesures de
\textit{Gyro\_X} :

\begin{equation}
    \hat{b}_\omega = \frac{1}{N} \sum_{i=1}^{N} \omega_{\text{mesuré},i}
\end{equation}

Ce biais est ensuite soustrait à l'ensemble des mesures de l'acquisition :

\begin{equation}
    \omega_{\text{corrigé}} = \omega_{\text{mesuré}} - \hat{b}_\omega
\end{equation}

\textbf{Justification du choix de l'axe :} comme observé lors de la mise en
place du matériel, l'accéléromètre mesure systématiquement $\approx 1g$ sur
l'axe X, confirmant que cet axe est orienté verticalement. La rotation du
piéton sur lui-même correspond donc à une rotation autour de l'axe vertical,
directement lue par \textit{Gyro\_X}.

\paragraph{Pour l'écart-type de l'accélération $\sigma_a$ pour la matrice $Q$}

La matrice de bruit de processus $Q$ modélise l'incertitude du modèle
cinématique. Elle dépend de $\sigma_a$, l'écart-type de l'accélération
horizontale du piéton. Ce paramètre ne peut pas être fixé arbitrairement :
il doit être estimé à partir des données réelles d'acquisition.

\textbf{Méthode RMS :} $\sigma_a$ est calculé comme la racine carrée de la
moyenne quadratique (RMS) des accélérations horizontales mesurées par l'IMU
(axes Y et Z, qui sont les axes horizontaux dans notre montage) :

\begin{equation}
    \sigma_a = \sqrt{\frac{1}{N} \sum_{i=1}^{N} \left( a_{Y,i}^2 + a_{Z,i}^2 \right)}
\end{equation}

Le choix du RMS plutôt que de la moyenne arithmétique est délibéré : la
moyenne des accélérations lors de la marche tend vers zéro par symétrie
(chaque pas produit une accélération suivie d'une décélération), ce qui
sous-estimerait l'énergie réelle du signal. Le RMS, en sommant les carrés,
capture correctement l'amplitude des pics d'accélération caractéristiques
du pas humain.

La valeur obtenue sur l'acquisition de référence est :

\begin{equation}
    \sigma_a = 0.86 \ \text{m/s}^2
\end{equation}

Ce paramètre est recalculé dynamiquement à chaque itération du filtre, 
permettant à $Q$ de s'adapter aux variations d'allure du piéton (marche
rapide, arrêts, changements de direction).